\section{Related Work}
\label{sec:related}

\subsection{Deep Learning for Fault Diagnosis}

Deep learning methods have been extensively applied to rotating machinery fault diagnosis, primarily through supervised learning on vibration data~\cite{wang2019deep,zhang2020convolutional}. CNN-based approaches process either raw temporal waveforms using 1D convolutions~\cite{inoue2018cnn} or time-frequency representations (spectrograms, scalograms) using 2D convolutions~\cite{wen2018new}. These methods learn hierarchical feature representations automatically, outperforming traditional feature engineering approaches that rely on hand-crafted statistical descriptors. More recently, Transformer-based architectures have been introduced to capture long-range dependencies in vibration signals~\cite{ding2022transformer}, leveraging self-attention to model global context. However, single-domain methods -- whether temporal-only, spectral-only, or purely Transformer-based -- fail to fully exploit the complementary information present in the time-frequency structure of vibration signals. Temporal waveforms capture transient impact dynamics with high resolution, while spectral representations reveal frequency content that characterizes specific fault modes. The inability to jointly model both domains limits diagnostic accuracy, particularly when signal characteristics shift across operating conditions.

\subsection{Few-Shot Learning for Fault Diagnosis}

The scarcity of labeled fault samples has motivated the application of few-shot learning (FSL) techniques to fault diagnosis. Meta-learning approaches, particularly optimization-based methods such as MAML~\cite{finn2017maml} and metric-based methods such as Prototypical Networks~\cite{snell2017prototypical} and Matching Networks~\cite{vinyals2016matching}, have been adapted for few-shot fault classification~\cite{li2021meta,sun2022multi}. These methods learn transferable knowledge across a distribution of tasks, enabling rapid adaptation to new fault types or operating conditions from limited examples. However, existing FSL frameworks for fault diagnosis face two key limitations. First, the meta-learning objective does not incorporate any physical understanding of the fault mechanisms, making the learned representations vulnerable to domain shifts that violate the statistical assumptions of the training distribution. Second, in extreme low-data regimes (1--5 shots per class), overfitting to spurious correlations in the limited support set degrades generalization performance. These limitations underscore the need for physically grounded regularization mechanisms in few-shot frameworks.

\subsection{Physics-Informed Fault Diagnosis}

Integrating physical domain knowledge into deep learning models has gained increasing attention in fault diagnosis research~\cite{wang2020physics}. Physics-Informed Neural Networks (PINNs) incorporate governing equations as soft constraints in the training objective~\cite{raissi2019physics}, while physics-guided architectures embed known signal properties into the network design. For example, wavelet-based convolutional layers replace standard convolutions with parameterized wavelet transforms to match the spectral characteristics of vibration signals~\cite{shen2023physics}. Dimensionless statistical metrics, such as kurtosis and skewness, have long been used in classical vibration analysis as indicators of fault severity~\cite{heng2009rotating}. These metrics are based on the physical insight that fault-induced impacts alter the probability distribution of vibration amplitudes in characteristic ways. Despite their diagnostic value, these physical priors are typically used as post-hoc features rather than being integrated into the learning process as inductive biases. Furthermore, existing physics-informed methods for fault diagnosis do not address the few-shot learning challenge, limiting their applicability in data-scarce industrial scenarios.

\subsection{Dual-Stream and Multi-Modal Architectures}

Several recent works have explored dual-stream architectures that combine temporal and spectral representations for vibration analysis~\cite{zhao2021fdfnet,han2022mcformer}. These methods typically employ two parallel encoders followed by a fusion mechanism. The key design choice lies in the fusion strategy: early fusion (concatenating features before classification), intermediate fusion (cross-modal interactions at intermediate layers), and late fusion (combining separate predictions). Existing approaches predominantly use simple fusion strategies such as feature concatenation, element-wise addition, or weighted averaging. While these approaches achieve improved performance over single-domain methods, they lack physically meaningful alignment between the temporal and spectral modalities. The fusion process treats both modalities as equally important without considering that their relative diagnostic value depends on the fault type and operating condition. A more principled fusion mechanism should align cross-modal representations using physically meaningful anchors that remain invariant under domain shifts.
